\begin{theorem}[固定尺度容量曲线的有限完备性]\label{thm:conclusion-capacity-finite-completeness}
固定分辨率 \(m\)。对折叠多重度 \(d_m(x)\) 定义
\[
h_d^{(m)}:=\#\{x\in X_m:\ d_m(x)=d\}
\qquad (1\le d\le 2^m),
\]
\[
N_m(T):=\#\{x\in X_m:\ d_m(x)\ge T\},
\]
\[
\mathcal C_m^{\mathrm{cont}}(T):=\sum_{x\in X_m}\min(d_m(x),T),
\]
\[
S_q(m):=\sum_{x\in X_m}d_m(x)^q.
\]
则下列四类数据在固定 \(m\) 下彼此唯一决定：
\begin{enumerate}
  \item 直方图 \(\{h_d^{(m)}\}_{d=1}^{2^m}\)；
  \item 尾计数函数 \(N_m\)；
  \item 连续容量曲线 \(\mathcal C_m^{\mathrm{cont}}\)；
  \item 有限矩列
  \[
  S_0(m),S_1(m),\dots,S_{2^m-1}(m).
  \]
\end{enumerate}
特别地，\(\mathcal C_m^{\mathrm{cont}}\) 不是粗摘要，而是纤维多重度多重集的一个完备编码。
\end{theorem}
\begin{proof}
由定理 \ref{thm:conclusion-wedderburn-threshold-holography}，数据 \(\textup{(1)}\)、\(\textup{(2)}\)、\(\textup{(3)}\) 在固定尺度 \(m\) 下已经彼此唯一决定；其中
\[
h_d^{(m)}=N_m(d)-N_m(d+1),
\qquad
\mathcal C_m^{\mathrm{cont}}(T)=\int_0^T N_m(t)\,\dd t.
\]
故只需证明 \(\textup{(1)}\iff\textup{(4)}\)。

由定义，
\[
S_q(m)=\sum_{d=1}^{2^m}h_d^{(m)}d^q
\qquad (q\ge 0).
\]
因此直方图显然决定有限矩列。

反过来，记
\[
h^{(m)}:=\bigl(h_1^{(m)},\dots,h_{2^m}^{(m)}\bigr)^\top,
\qquad
s^{(m)}:=\bigl(S_0(m),\dots,S_{2^m-1}(m)\bigr)^\top.
\]
则有线性系统
\[
s^{(m)}=V_m h^{(m)},
\qquad
(V_m)_{qd}=d^q,
\]
其中 \(0\le q\le 2^m-1\)、\(1\le d\le 2^m\)。这是节点
\[
1,2,\dots,2^m
\]
上的 Vandermonde 矩阵，因节点两两不同而可逆。故 \(h^{(m)}\) 由 \((S_q(m))_{q=0}^{2^m-1}\) 唯一恢复，从而 \(\textup{(4)}\Rightarrow\textup{(1)}\)。证毕。
\end{proof}
\leanverified{paper\_conclusion\_capacity\_finite\_completeness}

\begin{theorem}[容量--矩谱--零温端斜率的完备传输]\label{thm:conclusion-capacity-moment-zero-temperature-complete-transfer}
\leanverified{paper\_conclusion\_capacity\_moment\_zero\_temperature\_complete\_transfer}
设对每个整数 \(q\ge 1\)，极限
\[
\log r_q:=\lim_{m\to\infty}\frac{1}{m}\log S_q(m)
\]
存在；并设最大纤维指数极限
\[
\alpha_\ast:=\lim_{m\to\infty}\frac{1}{m}\log M_m,
\qquad
M_m:=\max_{x\in X_m}d_m(x),
\]
存在。则容量曲线族 \(\{\mathcal C_m^{\mathrm{cont}}\}_{m\ge 1}\) 已足以唯一决定整条高阶热力链
\[
\{r_q\}_{q\ge 1}
\quad\text{以及}\quad
\alpha_\ast.
\]
更精确地，
\[
S_q(m)=q\int_0^\infty T^{q-1}N_m(T)\,\dd T,
\]
从而
\[
\log r_q
=
\lim_{m\to\infty}
\frac{1}{m}\log\!\left(
q\int_0^\infty T^{q-1}N_m(T)\,\dd T
\right),
\]
并且
\[
\alpha_\ast=\lim_{q\to\infty}\frac{1}{q}\log r_q.
\]
因此，零温端斜率不需要额外的频谱输入；它是容量尾函数经 Mellin 变换后再取 \(q\to\infty\) 的纯尾部不变量。
\end{theorem}
\begin{proof}
由定理 \ref{thm:conclusion-capacity-finite-completeness}，整族容量曲线 \(\mathcal C_m^{\mathrm{cont}}\) 唯一决定全部尾计数函数 \(N_m\)。再由定理 \ref{thm:xi-capacity-moment-mellin-inversion-tail-capacity} 的 Mellin--Stieltjes 公式，
\[
S_q(m)=q\int_0^\infty T^{q-1}N_m(T)\,\dd T
\qquad (q>0).
\]
故每个 \(S_q(m)\) 以及相应极限 \(r_q\) 都由容量族唯一确定。

另一方面，结论 \ref{con:xi-fold-high-q-perron-endpoint-slope-maxfiber} 已经证明：若 \(\alpha_\ast\) 存在，则
\[
\lim_{q\to\infty}\frac{1}{q}\log r_q=\alpha_\ast.
\]
因此 \(\alpha_\ast\) 同样由 \(\{r_q\}_{q\ge 1}\) 唯一恢复，也就由容量曲线族唯一恢复。证毕。
\end{proof}
